\documentclass[a4,11pt]{aleph-notas}
% Se puede ver la documentación aquí:
% https://github.com/alephsub0/LaTeX_aleph-notas

% -- Paquetes adicionales
\usepackage{enumitem}
\usepackage{url}
\usepackage{array}
\usepackage{booktabs}
\usepackage{longtable}
\usepackage{aleph-comandos}
\hypersetup{
    urlcolor=blue,
    linkcolor=blue,
}


% -- Datos
\institucion{Facultad de Ciencias Exactas, Naturales y Ambientales}
\carrera{Catálogo STEM}
\asignatura{Matemática III}
\tema{Plan tema 15: Integrales de línea}
\autor{Andrés Merino}
\fecha{Periodo 2026-1}

\logouno[0.16\textwidth]{Logos/logoPUCE_04_ac}
\definecolor{colortext}{HTML}{1957d1}
\definecolor{colordef}{HTML}{1957d1}
\fuente{montserrat}


\begin{document}

\encabezado

%%%%%%%%%%%%%%%%%%%%%%%%%%%%%%%%%%%%%%%%
\section*{Resultado de Aprendizaje}
%%%%%%%%%%%%%%%%%%%%%%%%%%%%%%%%%%%%%%%%

%%%%%%%%%%%%%%%%%%%%%%%%%%%%%%%%%%%%%%%%
\subsection*{RdA de la asignatura:}
\begin{itemize}[leftmargin=*]
    \item \textbf{RdA 1:} Identificar conceptos, teoremas y propiedades de funciones vectoriales, derivadas parciales e integral vectorial.
    \item \textbf{RdA 2:} Calcular derivadas parciales e integrales con asistentes computacionales.
    \item \textbf{RdA 3:} Aplicar Cálculo Vectorial para resolver problemas reales.
\end{itemize}

%%%%%%%%%%%%%%%%%%%%%%%%%%%%%%%%%%%%%%%%
\subsection*{RdA del tema:}
\begin{itemize}[leftmargin=*]
    \item Calcular integrales de línea de campos escalares a lo largo de una curva parametrizada.
    \item Calcular integrales de línea de campos vectoriales y reconocer el efecto de la orientación de la curva.
    \item Emplear integrales de línea para determinar masa, centro de masas y momentos de inercia de alambres.
    \item Aplicar los teoremas fundamentales de las integrales de línea para campos conservativos.
\end{itemize}

%%%%%%%%%%%%%%%%%%%%%%%%%%%%%%%%%%%%%%%%
\section*{Introducción}
%%%%%%%%%%%%%%%%%%%%%%%%%%%%%%%%%%%%%%%%

\paragraph{Pregunta inicial:}
Cuando empujamos un carrito por un sendero que sube, baja y se curva, el esfuerzo total depende del camino que sigamos y de cuánto nos cueste avanzar en cada tramo. ¿Cómo podríamos sumar todas esas pequeñas contribuciones a lo largo de una trayectoria curva para obtener el esfuerzo total?

%%%%%%%%%%%%%%%%%%%%%%%%%%%%%%%%%%%%%%%%
\section*{Desarrollo}
%%%%%%%%%%%%%%%%%%%%%%%%%%%%%%%%%%%%%%%%

%%%%%%%%%%%%%%%%%%%%%%%%%%%%%%%%%%%%%%%%
\subsection*{Actividad 1: Integrales de línea de campos escalares y vectoriales}
Se introduce la integral de línea como un límite de sumas a lo largo de una curva, primero para campos escalares y luego para campos vectoriales, junto con sus aplicaciones físicas.

\paragraph{¿Cómo lo haremos?}
\begin{itemize}[leftmargin=*]
    \item \textbf{Motivación:} retomar la pregunta inicial e interpretar la integral de línea de un campo escalar como la masa de un alambre de densidad variable y la de un campo vectorial como el trabajo de una fuerza a lo largo de un camino.
    \item \textbf{Clase magistral:} se define la integral de línea de un campo escalar y se obtiene la fórmula $\int_C f\,d\alpha=\int_a^b f(\alpha(t))\|\alpha'(t)\|\,dt$; se define la integral de línea de un campo vectorial con la fórmula $\int_C F\cdot d\alpha=\int_a^b F(\alpha(t))\cdot\alpha'(t)\,dt$; se analiza el efecto de la orientación y de las parametrizaciones equivalentes. Se usa el resumen \href{https://andres-merino.github.io/MatematicaIII-05-N0051/01Resumen/Resumen15.pdf}{Resumen15.pdf}.
    \item \textbf{Implementación computacional:} calcular integrales de línea y visualizar curvas y campos con un asistente computacional.
    \item \textbf{Resolución de ejercicios:} calcular integrales de línea de campos escalares y vectoriales usando \href{https://andres-merino.github.io/MatematicaIII-05-N0051/04Ejercicios/Ejercicios15.pdf}{Ejercicios15.pdf}.
\end{itemize}

\paragraph{Verificación de aprendizaje:}
\begin{itemize}[leftmargin=*]
    \item Calcule $\int_C (x+y)\,ds$ donde $C$ es el segmento de $(0,0)$ a $(1,2)$.
    \item Calcule el trabajo del campo $F(x,y)=(y,x)$ a lo largo de la curva $\alpha(t)=(\cos t,\sin t)$, $t\in[0,\pi]$.
    \item ¿Qué ocurre con la integral de línea de un campo vectorial si se invierte la orientación de la curva? ¿Y con la de un campo escalar?
\end{itemize}

%%%%%%%%%%%%%%%%%%%%%%%%%%%%%%%%%%%%%%%%
\subsection*{Actividad 2: Teoremas fundamentales e independencia del camino}
Se estudia cuándo la integral de línea no depende del camino y se conecta con los campos conservativos a través de los teoremas fundamentales.

\paragraph{¿Cómo lo haremos?}
\begin{itemize}[leftmargin=*]
    \item \textbf{Motivación:} mostrar con un ejemplo que para ciertos campos el trabajo entre dos puntos es el mismo sin importar el camino, y preguntar qué caracteriza a esos campos.
    \item \textbf{Clase magistral:} se define la independencia del camino; se presenta la equivalencia entre campo conservativo, integral independiente del camino e integral nula sobre curvas cerradas; se enuncian el Primer y Segundo Teorema Fundamental para integrales de línea. Se usa el resumen \href{https://andres-merino.github.io/MatematicaIII-05-N0051/01Resumen/Resumen15.pdf}{Resumen15.pdf}.
    \item \textbf{Resolución de ejercicios:} evaluar integrales de línea de campos conservativos mediante una función potencial usando \href{https://andres-merino.github.io/MatematicaIII-05-N0051/04Ejercicios/Ejercicios15.pdf}{Ejercicios15.pdf}.
\end{itemize}

\paragraph{Verificación de aprendizaje:}
\begin{itemize}[leftmargin=*]
    \item Dado el campo conservativo $F(x,y)=(2x,2y)$ con potencial $f(x,y)=x^2+y^2$, calcule $\int_C F\cdot ds$ de $(0,0)$ a $(1,1)$.
    \item ¿Por qué la integral de un campo conservativo sobre una curva cerrada es igual a cero?
    \item Enuncie las tres condiciones equivalentes para que un campo vectorial sea conservativo.
\end{itemize}

%%%%%%%%%%%%%%%%%%%%%%%%%%%%%%%%%%%%%%%%
\section*{Cierre}
%%%%%%%%%%%%%%%%%%%%%%%%%%%%%%%%%%%%%%%%

\paragraph{Tarea:}
Resolver del libro \href{https://catalogobiblioteca.puce.edu.ec/cgi-bin/koha/opac-detail.pl?biblionumber=83405&query_desc=su%3A%22An%C3%A1lisis%20vectorial%22}{Cálculo Vectorial de Colley}: sección 6.1, ejercicios 1--13 (impares).

\paragraph{Pregunta de investigación:}
\begin{enumerate}[leftmargin=*]
    \item ¿Cómo se interpreta físicamente la integral de línea de un campo vectorial como trabajo? Investigar un ejemplo concreto en mecánica.
    \item ¿Qué significa que un campo de fuerzas sea conservativo en física? ¿Qué relación tiene con la energía potencial?
\end{enumerate}

\paragraph{Para el próximo tema:}
Ver los videos:
\begin{itemize}[leftmargin=*]
    \item \href{https://youtu.be/_GHJB0qbp_E?si=eY4iZ8Yk2hX6m3vO}{Introduction to the surface integral}.
\end{itemize}


\end{document}
